%%-*-latex-*-

\documentclass[wide]{slides}

% Colour in verbatim
%
\usepackage{alltt}
\usepackage{xcolor}

% Language
%
\usepackage{hyphenat}          % \hyp{} is a breakable dash
\usepackage{alltt}

\frenchspacing  % Follow French conventions after a period

% Maths & Logic
%
\usepackage{stmaryrd}
\usepackage{mathpartir}

% URLs
\usepackage{url}

% Graphics
%
\usepackage{graphicx}

% Code listing
%
\usepackage{listings}

\usepackage{color}
\definecolor{lightgray}{rgb}{.9,.9,.9}
\definecolor{darkgray}{rgb}{.4,.4,.4}
\definecolor{purple}{rgb}{0.65, 0.12, 0.82}

\lstdefinelanguage{JavaScript}{
  keywords={match, when, type, const, let, typeof, new, true, false, catch, function, return, null, catch, switch, var, if, in, while, do, else, case, break},
  keywordstyle=\color{blue}\bfseries,
  ndkeywords={class, export, boolean, throw, implements, import, this},
  ndkeywordstyle=\color{darkgray}\bfseries,
  identifierstyle=\color{black},
  sensitive=false,
  comment=[l]{//},
  morecomment=[s]{/*}{*/},
  commentstyle=\color{purple}\ttfamily,
  stringstyle=\color{red}\ttfamily,
  morestring=[b]',
  morestring=[b]"
}

\lstset{
   language=JavaScript,
   backgroundcolor=\color{lightgray},
   extendedchars=true,
   basicstyle=\footnotesize\ttfamily,
   showstringspaces=false,
   showspaces=false,
   numbers=left,
   numberstyle=\footnotesize,
   numbersep=9pt,
   tabsize=2,
   breaklines=true,
   showtabs=false,
   captionpos=b
}

% New environments and commands
%
\newcommand{\JsLIGO}[0]{\textsf{JsLIGO}}
\newcommand{\CameLIGO}[0]{\textsf{CameLIGO}}
\newcommand{\OCaml}[0]{\textsf{OCaml}}
\newcommand{\Menhir}[0]{\textsf{Menhir}}
\newcommand{\TypeScript}[0]{\textsf{TypeScript}}
\newcommand{\TreeSitter}[0]{\textsf{tree-sitter}}
\newcommand{\Michelson}[0]{\textsf{Michelson}}

% Input files
%
% \input{commands}

% ----------------------------------------------------------------
% Document
%
\maintitle{LIGO for TypeScript developers on Layer-1}
\mainauthor{\textbf{Christian Rinderknecht}\\
  {\small\url{Christian.Rinderknecht@trili.tech}}\\
\textsf{TriliTech}}
\confname{Show and Tell}
\confshortname{S\&T}
\confdate{17 April 2025}

\begin{document}

\maketitle

\begin{slide}
  \title{A short history of LIGO}

  \begin{itemize}

    \item I did my PhD at INRIA under the supervision of Michel Mauny.

    \item I joined Nomadic in July 2018, where I started designing a
      smart contract language, similar to \textsf{Pascal}, that would
      be compiled to \Michelson{}.

    \item My then-colleague Gabriel Alfour also wanted to make
      programming smart contracts on \textsf{Tezos} easier, but with a
      bottom\hyp{}up approach: building a set of macros that expand
      into \Michelson{}.

    \item We joined our efforts in 2019 into what became LIGO.

    \item I contributed the design and the front-end, he did the
      type-checker and code generator (50-50 in lines of code).

  \end{itemize}
\end{slide}

\begin{slide}
  \title{A short history of LIGO}

  \begin{itemize}

    \item The idea was for LIGO to offer different concrete syntaxes
      inspired from mainstream programming languages.

    \item Those LIGO languages were meant to be domain-specific, but
      not overly so: supporting special literals (natural numbers,
      mutez, for examples), but taking most Tezos-specific types and
      functions from a library.

    \item The front-end had to lower their input to the common
      pipeline, while sharing as much code as possible in the
      processing of the different syntaxes.

    \item The first syntax was \textsf{PascaLIGO}, as I thought it
      would suit an imperative style, and lots of keywords would make
      the code more explicit, given the high stakes.

    \item Next, turmoil in the nascent ecosystem lead me to design
      \CameLIGO{}. Then came \textsf{ReasonLIGO},
      \JsLIGO{} (which is today's topic), and even
      \textsf{PyLIGO}, which was never plugged in.

  \end{itemize}
\end{slide}

\begin{slide}
  \title{Architecture}

  \begin{itemize}

    \item Nowadays, only two LIGO syntaxes remain: \CameLIGO{}
      and \JsLIGO{}. The former is a strict subset of
      \OCaml{}, if we leave aside the Tezos-specific literals, a
      local type definition construct, and attributes (they occur in
      prefix position in LIGO).

    \item \JsLIGO{} started also as a subset but evolved and
      became syntactically incompatible with its model,
      \TypeScript{}. For example, a proposal for pattern
      matching in \TypeScript{} was implemented in
      \JsLIGO{}.

    \item The front-end is made, in order, of a preprocessor, a lexer,
      two parsers and the first tree transformation before
      type\hyp{}inference. Each pass can be built as a standalone
      executable for testing.

    \item I wrote the preprocessor because, at the time, we had no
      time to design and implement a module system, together with a
      build system.

    \item Unfortunately, it became very popular, even when LIGO got
      itself a module system, because it enabled conditional
      compilation (often for new type definitions of the storage,
      depending on the version of the standard being implemented).

  \end{itemize}
\end{slide}

\begin{slide}
  \title{Preprocessor and Lexer}

  \begin{itemize}

    \item The preprocessor is not trivial, as it needs to understand
      UTF-8 encoded glyphs (in comments only) and the kind of comments
      and strings that depend on the LIGO variant.

    \item The lexer is also a functor parameterised by a lexer that is
      specific to a given LIGO syntax (comments, strings, keywords
      etc.). Error reporting has to be aware of linemarkers left by
      the preprocessor. Some style constraints are enforced (like a
      linter) for the sake of clarity.

    \item The lexer is specified as an input to \textsf{ocamllex}.

    \item Until \JsLIGO{} was added to the party, the lexer
      signature was a function that returned a single token, which the
      parser called on demand. Syntactic constructs difficult to parse
      sometimes rely on lexer hacks, that is, having the parser
      informing the lexer of the syntactic left\hyp{}context. Instead,
      I had the lexer run on all the (preprocessed) input, before
      applying a series of self\hyp{}passes on the tokens
      ---~injecting virtual tokens when needed to help the
      parser. This is not as expressive as having the syntactic
      context, but rational approximations work quite well in
      practice.

  \end{itemize}
\end{slide}

\begin{slide}
  \lstset{language=JavaScript}

  \title{Lexer}

  \begin{itemize}

    \item Perhaps the most interesting self\hyp{}pass on the tokens is
      found with \JsLIGO{}:
      \begin{lstlisting}
        const f = ([x,y]) ...
      \end{lstlisting}
      could be a parenthesised pair or the start of an arrow function
      (a shift\hyp{}reduce conflict). Analysing the lexical
      right\hyp{}context in the grammar helps disambiguate the two
      cases, and the lexer injects a \verb+ES6FUN+ token in front of a
      function:
      \begin{quote}
      \verb|CONST IDENT EQ ES6FUN LPAR ...|
      \end{quote}
      The parser knows a function is coming next upon reading the
      virtual token \verb+ES6FUN+ (no more conflict).

    \item A classic one is needed to handle ``\verb+>>+'' in
      \begin{lstlisting}
        type t = u<v<w>>
      \end{lstlisting}
      for which a self\hyp{}pass injects a Zero\hyp{}Width SPace
      (\verb|ZWSP|) virtual token between the two closing chevrons:
      \begin{quote}
       \verb|TYPE IDENT EQ IDENT LCHEV IDENT LCHEV IDENT RCHEV ZWSP RCHEV|
      \end{quote}
      The parser can then distinguish the binary operator
      ``\verb+>>+'' from the above.

  \end{itemize}
\end{slide}

\begin{slide}
  \title{The parser generator Menhir}

  \begin{itemize}

    \item The LIGO pipeline is more than a compiler, because it
      provides support for a LSP, and that is why the front-end has to
      keep all the source information, including comments. This adds
      significant complexity to the lexer, which has to hide it from
      the parser.

    \item The two LIGO parsers are generated by \Menhir{}.

    \item \Menhir{} accepts LR(1) grammars, with semantic
      actions in \OCaml{}. In practice, many grammars are
      annotated with precedence declarations to solve
      LR\hyp{}conflicts, potentially making the grammars not LR(1),
      and making it difficult to reason about them.

    \item Both \CameLIGO{} and \JsLIGO{} are given pure
      LR(1) grammars, thanks to the rewriting of productions, and the
      use of virtual tokens.

    \item Yann Régis-Gianas wrote the front-end of Menhir, and
      François Pottier the automaton construction and the
      back\hyp{}end. The original interest was an application of GADTs
      to the static typing of the parse stack: \url{https://cambium.inria.fr/~fpottier/publis/fpottier-regis-gianas-typed-lr.pdf}
  \end{itemize}
\end{slide}

\begin{slide}
  \title{Menhir and the JsLIGO Parser}

  \begin{itemize}

    \item \Menhir{} can identify all error states in the
      LR\hyp{}automaton, so we can write precise error messages for
      them. They are precise because the stack and syntactic
      righ\hyp{}context is available at the error state, instead of
      just only the next valid tokens, so the past (the intention of
      the programmer, inferred from the prefix of the stack in terms
      of grammatical productions) can be expressed in terms of
      constructs, and the future too, if we take some care when
      writing the
      grammar. \url{https://cambium.inria.fr/~fpottier/publis/bour-pottier-reachability.pdf}

    \item The fun part of the \JsLIGO{} grammar is the handling
      of the (in)famous \emph{Automatic Semicolon Insertion
      (ASI)}. JavaScript has a set of dynamic heuristics allowing in
      some contexts to omit a semicolon. Instead, I analysed the
      grammar and hardcoded a sub\hyp{}grammar (\~{}100~LOC) for
      statements that allow some semicolons to be missing (and omits
      more than the mainstream compilers!).

    \item One last very cool paper about using \Menhir{} to
      parse C11 (as in \textsf{CompCert}):
      \url{https://cambium.inria.fr/~fpottier/publis/jourdan-fpottier-2016.pdf}
  \end{itemize}
\end{slide}

\begin{slide}
  \title{Error recovery with Menhir}

  \begin{itemize}

    \item \Menhir{} has several interfaces, improperly called
      APIs. The \emph{Incremental API} enables the client to take
      control when an error is encountered.

    \item The parse stack (of terminal and non\hyp{}terminals
      recognised so far) and the automaton state (e.g. the
      LR\hyp{}items) are available for inspection, as persistent
      values.

    \item This enables to decide to either handle the error or try and
      recover from it: a strategy has to be provided by the client.

    \item LIGO features error recovery for both \CameLIGO{} and
      \JsLIGO{}. Two continuations are possible for recovery: either
      skipping some tokens or inserting one. The decision is made
      based on a heuristic derived from a study of the grammar and
      common cases (design patterns, keywords and symbol
      uses,tabulations etc.).

    \item The LSP makes constant use of this feature to maintain
      syntax high\hyp{}lighting, to display types on hover, and
      for go\hyp{}to\hyp{}definition on contracts being edited.

  \end{itemize}
\end{slide}


\begin{slide}
  \title{The Unified AST}

  \begin{itemize}

    \item The parser produces an Abstract Syntax Tree (AST)
      ---~confusingly named Concrete Syntax Tree (CST) in the code
      base for legacy reasons: it used to be what its name proclaims.

    \item The AST is then transformed into a more abstract kind of
      tree, the \emph{unified AST}. Whereas the CSTs are specific to
      each LIGO variant, the unified AST is common to all, and can be
      considered as the entrance to the common pipeline.

    \item Because it is the target of a functional and imperative
      language, respectively \CameLIGO{} and \JsLIGO{},
      the definition of the Unified AST is cumbersome. Also, in order
      to enable some measure of decompilation (type expressions for
      the LSP), it cannot be too abstracted.

  \end{itemize}
\end{slide}

\begin{slide}
  \title{JsLIGO concrete syntax}

\begin{quote}
\begin{lstlisting}
type parameter = unit;
type storage = unit;
type t = [list<operation>, storage];

@entry
function give5tez (param: parameter, storage: storage): t {
  let operations: list<operation> = [];
  if (Tezos.get_balance() >= 5tez) {
    const receiver =
      match(Tezos.get_contract_opt(Tezos.get_sender())) {
        when(Some(contract)): contract;
        when(None): failwith("Couldn't find account");
      };
    operations =
      [Tezos.Operation.transaction(param, 5tez, receiver)];
  }
  return [operations, storage];
}
\end{lstlisting}
\end{quote}

\end{slide}

\begin{slide}
  \title{JsLIGO versus TypeScript}

Even though the syntax error messages displayed by the TypeScript
Playground, for example, are relatively good, the \JsLIGO{} parser
enables to be correct about the past and complete about the
future. For instance:
\begin{quote}
\small
\begin{alltt}
File "syn_err.jsligo", line 2, characters 9-11:
  1 | interface I \{
  2 |   type t \textcolor{red}{+=} int;
  3 | \}
\textcolor{red}{Ill-formed type declaration in an interface.
At this point, one of the following is expected:
  * the assignment symbol '=' followed by a type;
  * type parameters (variables) between chevrons ('<', '>') if the
    type is generic;
  * a semicolon ';' if the type is abstract;
  * a closing brace '\}' if no more entries.}
\end{alltt}
\end{quote}

\end{slide}

\begin{slide}
  \title{JsLIGO versus TypeScript}

  \begin{itemize}

    \item Despite the subset of TypeScript needed for writing
      \JsLIGO{} contracts is comfortable enough (some computations can
      be expressed in several ways, for instance, and insertion of
      \Michelson{} code is possible), a syntax error cannot
      tell whether the input is actually valid \TypeScript{} or not.

    \item For example, in the previous example, the parser cannot
      actually tell that the input is not valid \TypeScript{}, and the
      penultimate example is valid \JsLIGO{} but not correct
      \TypeScript{}: \JsLIGO{} is a non\hyp{}conservative extension of
      \TypeScript{}.

    \item Likewise, a type error cannot tell whether the source is
      valid \TypeScript{}. Worse, since the LIGO type system is
      \emph{ad hoc}, the inferred types may differ from those of a
      \TypeScript{} compiler (for example, due to different implicit
      type coercions).

    \item The LIGO tooling (mainly the LSP) is custom\hyp{}made.

  \end{itemize}
In short: Idioms and design patterns commonly used by \TypeScript{}
developers are not supported in \JsLIGO{}; it is not possible for the
LIGO compiler to tell whether it is because the input is not valid
\TypeScript{}. If the code compiles, but the semantics differs, the
LIGO team would consider it a bug, but it could be silent (we never
found an example, though).

\end{slide}

\begin{slide}
  \title{The parser generator tree-sitter}

  \begin{itemize}

    \item In order to leverage existing \TypeScript{} tooling, for
      instance, a pretty\hyp{}printer,

    \item to have the LIGO compiler inform the developers of the
      differences between \TypeScript{} and \JsLIGO{},

    \item and to avoid writing a lexer and parser for \TypeScript{}
      with \textsf{ocamllex} and \Menhir{}, a new approach was
      proposed, namely use a third\hyp{}party \TypeScript{} parser.

    \item Since those are most often closed\hyp{}source or blatantly
      incomplete, we decided to use a parser in C generated by
      \TreeSitter{} from a \TypeScript{} grammar.

    \item It is difficult to assess the correctness, let alone the
      conformance, of that parser.

    \item First, contrary to \Menhir{}, which has a \textsf{Coq}
      back\hyp{}end, we do not have a way to prove that the parser
      indeed recognises the language generated by the grammar, and is
      complete, that is, exactly the valid inputs are
      accepted. (Fun/Sad fact: termination is not proved with
      ``\texttt{menhir --coq}''!)

    \item Second, there is no more formal grammar of \TypeScript{}
      published by Microsoft since version 1.8 (current is 5.8.3)
      ---~in 2016.

  \end{itemize}

\end{slide}

\begin{slide}
  \title{The parser generator tree-sitter}

  \begin{itemize}

    \item \TreeSitter{} has been designed to generate parsers in C
      from LR(1) grammars.

    \item Those parsers perform \emph{error recovery}: when
      encountering an unexpected token, the parser can either
      \emph{skip} any number of tokens or non\hyp{}terminal symbols
      from the input, or \emph{insert} one virtual terminal (token)
      when one is deemed missing in the input.

    \item Instead of chosing one strategy rather than the other, the
      parser engine runs internally as a \emph{Generalised LR} (GLR)
      algorithm, so, in general, the parse stack is instead a tree of
      valid prefix sentences (a.k.a. sentential forms). Skipping and
      inserting become two new branches (up to six).

    \item Each branch of the tree maintains two quantities: an
      \emph{error cost} and a \emph{node count}, respectively, an
      integer function of the number of skipped items and inserted
      tokens, and the number of valid subtrees parsed since the last
      error.

    \item One heuristic prunes the automaton's tree based on those two
      values, and another selects one final parse tree if many valid
      were found (smaller and earlier trees are preferred).

  \end{itemize}

\end{slide}

\end{document}
